\section*{الفصل \ref{ch:g6:inside-the-earth} --- باطن الأرض}

\begin{solution}{exo:g6:inside-the-earth:1}
\omterm{def:g6:inside-the-earth:layers}{القشرة}، \omterm{def:g6:inside-the-earth:layers}{والوشاح}، \omterm{def:g6:inside-the-earth:layers}{واللب الخارجي}، \omterm{def:g6:inside-the-earth:layers}{واللب الداخلي}. \omterm{def:g6:inside-the-earth:layers}{اللب الخارجي}
سائل. \omterm{def:g6:inside-the-earth:layers}{واللب الخارجي} \omterm{def:g6:inside-the-earth:layers}{واللب الداخلي} مكوَّنان من الفلز (معظمه حديد، مع شيء من
النيكل)؛ أما \omterm{def:g6:inside-the-earth:layers}{القشرة} \omterm{def:g6:inside-the-earth:layers}{والوشاح} فمكوَّنان من \omterm{def:g3:stones-pebbles-sand:rock}{الصخر}.
\end{solution}

\begin{solution}{exo:g6:inside-the-earth:2}
(أ) خطأ: \omterm{def:g6:inside-the-earth:layers}{القشرة} أسمك تحت \omterm{def:g3:continents-and-oceans:continent}{القارات} (نحو 30~كيلومترًا) منها
تحت المحيطات (نحو 5~كيلومترات). (ب) خطأ: \omterm{def:g6:inside-the-earth:layers}{اللب} مكوَّن من الفلز، وهو
سبيكة من الحديد والنيكل. (ج) صواب. (د) صواب: 12\,262~مترًا، أي نحو 12.3~كيلومتر.
\end{solution}

\begin{solution}{exo:g6:inside-the-earth:3}
(أ) عند 2891~كيلومترًا، أي نحو 2900~كيلومتر. (ب) عند 5150~كيلومترًا. (ج) $5150 - 2891 = 2259$:
نحو 2260~كيلومترًا.
\end{solution}

\begin{solution}{exo:g6:inside-the-earth:4}
«القشر يشبه \underline{\omterm{def:g6:inside-the-earth:layers}{القشرة}}، والبياض يشبه
\underline{\omterm{def:g6:inside-the-earth:layers}{الوشاح}} والصفار يشبه \underline{\omterm{def:g6:inside-the-earth:layers}{اللب}}. لكن خلافًا
للصفار، يتكون \underline{\omterm{def:g6:inside-the-earth:layers}{اللب}} من جزأين: \underline{\omterm{def:g6:inside-the-earth:layers}{لب خارجي}}
سائل و\underline{\omterm{def:g6:inside-the-earth:layers}{لب داخلي}} صلب.»
\end{solution}

\begin{solution}{exo:g6:inside-the-earth:5}
$8.849 \div 6371 \approx 0.00139$، و$0.00139 \times 100 \approx 0.14$:
\omterm{def:g4:world-physical-map:altitude}{فارتفاع} إفرست لا يتجاوز نحو \textbf{0.14~\%} من نصف القطر. وإذا نظرنا من بعيد،
فالأرض كرة ملساء جدًا: وأعلى جبالها نتوء ضئيل.
\end{solution}

\begin{solution}{exo:g6:inside-the-earth:6}
(أ) $30 \div 5 = 6$: أسمك بست مرات. (ب) $100 \div 5 = 20$: أسمك بعشرين
مرة.
\end{solution}

\begin{solution}{exo:g6:inside-the-earth:7}
(أ) $12 + 1 \times 25 = 37~^\circ$C. (ب) $12 + 1.5 \times 25 = 12 + 37.5 =
49.5~^\circ$C. (ج) $12 + 2 \times 25 = 62~^\circ$C.
\end{solution}

\begin{solution}{exo:g6:inside-the-earth:8}
الارتفاع $74 - 14 = 60$ درجة على امتداد 2~كيلومتر، إذن $60 \div 2 =
\textbf{30}$~درجة لكل كيلومتر. نعم: 30 تقع بين متوسط الشرق
(25) ومتوسط الغرب (34).
\end{solution}

\begin{solution}{exo:g6:inside-the-earth:9}
(أ) فوق \omterm{def:g6:inside-the-earth:layers}{اللب} مباشرة، نحو 13.7~كيلومتر في الثانية؛ وتحت قمته مباشرة،
نحو 8.1~كيلومتر في الثانية: فالموجات P تتباطأ فجأة. (ب) بين
2891~كيلومترًا و5150~كيلومترًا، في \omterm{def:g6:inside-the-earth:layers}{اللب الخارجي}. (ج) الموجات S لا تستطيع أن تنتقل عبر
السائل: \omterm{def:g6:inside-the-earth:layers}{فاللب الخارجي} سائل.
\end{solution}

\begin{solution}{exo:g6:inside-the-earth:10}
(أ) $12.3 \div 30 = 0.41$: نحو \textbf{41~\%} من \omterm{def:g6:inside-the-earth:layers}{قشرة} قارية
متوسطة. (ب) لا: تحت \omterm{def:g3:continents-and-oceans:continent}{القارة} يبدأ \omterm{def:g6:inside-the-earth:layers}{الوشاح} تحت \omterm{def:g6:inside-the-earth:layers}{القشرة}، على عمق نحو
30~كيلومترًا في المتوسط؛ وقد توقف ثقب الحفر بعيدًا فوقه، داخل \omterm{def:g6:inside-the-earth:layers}{القشرة}.
\end{solution}

\begin{solution}{exo:g6:inside-the-earth:11}
(أ) $29 \times 6371 = 184\,759$، و$15 + 184\,759 = 184\,774$: نحو
185\,000~درجة. (ب) يعطي المختبر نحو 6000~درجة عند حافة
\omterm{def:g6:inside-the-earth:layers}{اللب الداخلي}: فعدد ليو أكبر من الحقيقة بنحو 30~مرة
($184\,774 \div 6000 \approx 31$). (ج) معدل 29~درجة لكل كيلومتر
لم يُقس إلا في أول كيلومترين تقريبًا. وليو يستعمله على بعد آلاف
الكيلومترات من القياسات؛ وفي الأعماق ترتفع درجة الحرارة ببطء
أكبر بكثير. ولا يجوز استعمال معدل إلا حيث قيس.
\end{solution}

\begin{solution}{exo:g6:inside-the-earth:12}
(أ) إذا كانت \omterm{def:g6:inside-the-earth:layers}{القشرة} أخف من المتوسط، فلا بد أن تكون المادة الموجودة في الأعماق
أكثف من 5.5، لترفع المتوسط. (ب) في \omterm{def:g6:inside-the-earth:layers}{اللب الخارجي} \omterm{def:g6:inside-the-earth:layers}{واللب
الداخلي} (من نحو 9.9 إلى نحو 13~غرامًا لكل سنتيمتر مكعب)؛ أما
\omterm{def:g6:inside-the-earth:layers}{الوشاح} فبالكاد يبلغ 5.6 في قاعه. (ج) المواد الأكثف تغوص
تحت المواد الأخف، كما يغوص الحجر في الماء: فانتهى الفلز الثقيل في
الوسط \omterm{def:g3:stones-pebbles-sand:rock}{والصخور} الأخف في الأعلى.
\end{solution}

\begin{solution}{pb:g6:inside-the-earth:1}
\textbf{1.} $2 \times 6371 = 12\,742$~كيلومترًا. وقطر النموذج 2~متر، أي
2000~مليمتر.

\textbf{2.} $12\,742 \div 2000 \approx 6.37$: يمثّل 1~مليمتر من النموذج
نحو 6.37~كيلومتر من الأرض الحقيقية.

\textbf{3.} $8.849 \div 6.371 \approx 1.39$: \omterm{def:g4:world-physical-map:altitude}{ارتفاع} إفرست على النموذج نحو
\textbf{1.4~مليمتر}، أي حبة \omterm{def:g3:stones-pebbles-sand:sand}{رمل}.

\textbf{4.} $12.3 \div 6.371 \approx 1.9$~mm. ومركز النموذج على بعد
1000~مليمتر (1~متر) تحت السطح: فثقب الحفر أقصر بنحو 520~مرة
($1000 \div 1.9 \approx 520$).

\textbf{5.} $2891 \div 6.371 \approx 453.8$: تُرسم قمة \omterm{def:g6:inside-the-earth:layers}{اللب الخارجي}
على بعد نحو \textbf{454~مليمترًا} تحت السطح.

\textbf{6.} $5150 \div 6.371 \approx 808.3$: على بعد نحو 808~مليمترات تحت
السطح. ونصف قطر \omterm{def:g6:inside-the-earth:layers}{اللب الداخلي} على النموذج $1000 - 808 = 192$~mm.

\textbf{7.} \omterm{def:g6:inside-the-earth:layers}{الوشاح}: من نحو 5~مليمترات إلى 454~مليمترًا، أي نحو 449~مليمترًا. ويشغل \omterm{def:g6:inside-the-earth:layers}{اللب
الخارجي} $808 - 454 = 354$~mm \omterm{def:g6:inside-the-earth:layers}{واللب الداخلي} 192~مليمترًا.

\textbf{8.} $30 \div 6.371 \approx 4.7$~mm: شريط رقيق، يمكن
رسمه بفرشاة دقيقة.

\textbf{9.} $15 + 2 \times 29 = 73~^\circ$C. وعلى النموذج، يقع عمق 2~كيلومتر على بعد
$2 \div 6.371 \approx 0.31$~mm تحت السطح: بالكاد سمك
ثلاث ورقات.

\textbf{10.} $4 \div 6.371 \approx 0.63$~mm، مقابل نحو 1.9~مليمتر لثقب
الحفر: فالمنجم أقل عمقًا، بنحو $1.9 - 0.6 = 1.3$~mm.

\textbf{11.} انصهار المادة يتوقف على درجة الحرارة وعلى
الضغط. \omterm{def:g6:inside-the-earth:layers}{واللب الداخلي} أشد الطبقات انضغاطًا، بنحو
3.3~مليون ضعف ضغط الهواء: والحديد المضغوط لا ينصهر إلا عند
درجة حرارة أعلى، لذلك يبقى حديد \omterm{def:g6:inside-the-earth:layers}{اللب الداخلي} صلبًا مع أنه
أحر من \omterm{def:g6:inside-the-earth:layers}{اللب الخارجي}.

\textbf{12.} $5 \div 6.371 \approx 0.785$: يجب ألا يزيد سمك \omterm{def:g6:inside-the-earth:layers}{القشرة} المحيطية المرسومة
على \textbf{0.78~مليمتر} في النموذج الذي قطره مترَان، أي سمك
نحو ثماني ورقات ($0.78 \div 0.1 \approx 8$).
\end{solution}
